\uuid{RmYS}
\titre{Logique}
\auteur{bodin ; croizat ; sacre}
\organisation{exo7}
\licence{CC BY-NC-SA 4.0}
\datecreate{2022-03-08}
\niveau{L1}
\module{Algèbre}
\chapitre{Logique, ensemble, raisonnement}
\sousChapitre{Logique}
\difficulte{3}
\theme{Logique}
\exo7qcmid{hilisit.logique.013}

\contenu{
\begin{question}
\qclassification{100.01}
$f : \Rr \to \Rr$ désigne une fonction.
Pour les phrases suivantes dire si la négation proposée est correcte.
\begin{answers}
	\good{La négation de "Il existe $x\in\Rr$ tel que $f(x)=0$" est "Pour tout $x\in\Rr$ on a $f(x)\neq0$".}
	\good{La négation de "Pour tout $x\in\Rr$ on a $f(x)=0$" est "Il existe $x\in\Rr$ tel que $f(x)\neq0$".}
	\good{La négation de "Il existe $x\in\Rr$ tel que $f(x)\ge0$" est "Pour tout $x\in\Rr$ on a $f(x)<0$".}
	\bad{La négation de "Pour tout $x\in\Rr$ on a $f(x)>0$" est "Pour tout $x\in\Rr$ on a $f(x)\le0$".}
\end{answers}
\begin{explanations}
La négation de "Pour tout" est "Il existe" (et réciproquement). La négation de "$f(x)=0$" est "$f(x) \neq 0$".  La négation de "$f(x)\ge0$" est "$f(x) < 0$".

Toutes les affirmations sont vraies sauf la négation de "Pour tout $x\in\Rr$ on a $f(x)>0$" qui est "Il existe $x\in\Rr$ tel que $f(x)\le0$".
\end{explanations}
\end{question}
}
